% Generated draft: requires mathematical and source review.
\documentclass[11pt]{article}
\usepackage[margin=25mm]{geometry}
\usepackage{fontspec}
\setmainfont{Latin Modern Roman}
\usepackage{amsmath,amssymb,mathtools,unicode-math}
\setmathfont{Latin Modern Math}
\usepackage{enumitem,needspace,xurl,hyperref}
\setlength{\parindent}{0pt}
\setlength{\parskip}{5pt}
\setlength{\emergencystretch}{3em}
\newcommand{\N}{\mathbb{N}}
\newcommand{\Z}{\mathbb{Z}}
\newcommand{\Q}{\mathbb{Q}}
\newcommand{\R}{\mathbb{R}}
\newcommand{\C}{\mathbb{C}}
\newcommand{\F}{\mathbb{F}}
\newcommand{\A}{\mathbb{A}}
\newcommand{\PP}{\mathbb{P}}
\newcommand{\E}{\mathbb{E}}
\newcommand{\eps}{\varepsilon}
\newcommand{\dd}{\,\mathrm{d}}
\begin{document}
\section*{UnsolvedMath ID: 701258}
\textbf{Source problem code:} catalogue entry
\subsection{Definitions and mathematical statement}
Let \(G\) be a group, with identity subgroup denoted by \(1\). Define its lower central series by \(\gamma_1(G)=G\) and \(\gamma_{k+1}(G)=[\gamma_k(G),G]\). The group has nilpotency class exactly \(n\geq 1\) if \(\gamma_{n+1}(G)=1\) and \(\gamma_n(G)\neq 1\). No finiteness or finite-generation hypothesis on \(G\) is imposed.

Define the upper central series by \(\zeta_0(G)=1\) and
\[
\zeta_{k+1}(G)/\zeta_k(G)=Z\bigl(G/\zeta_k(G)\bigr).
\]
The source calls \(D_k(G)\) the \(k\)-th dimension subgroup without specifying coefficients. Under the conventional integral interpretation, let \(\varepsilon:\Z[G]\to\Z\) be the augmentation map, \(\varepsilon(\sum_g a_g g)=\sum_g a_g\), let \(\Delta(G)=\ker\varepsilon\), and set
\[
D_k(G)=\{g\in G:g-1\in\Delta(G)^k\},\qquad k\geq 1.
\]
This coefficient convention requires confirmation from the original source; the formulation below is conditional on it.

For each integer \(n\geq 1\), put
\[
A_n=\{m\in\Z_{\geq 1}:\text{there exists a group }G\text{ of nilpotency class exactly }n\text{ with }D_m(G)\neq 1\},
\]
\[
B_n=\{m\in\Z_{\geq 0}:\text{there exists a group }G\text{ of nilpotency class exactly }n\text{ with }D_n(G)\not\subseteq\zeta_m(G)\}.
\]
The source writes
\[
f(n)=\max A_n,\qquad g(n)=\max B_n.
\]
Here the nonnegative index convention for \(B_n\) is explicit: the source does not specify whether \(m=0\) is allowed. Restricting to positive indices can leave this set empty, already when \(n=1\). The notation \(\max A_n\) does not assert that a finite maximum exists; part (b) asks precisely for boundedness of \(A_n\). The displayed definition of \(g\) retains \(D_n(G)\) exactly as supplied. Since \(\zeta_n(G)=G\) for a group of class \(n\), its admissible indices satisfy \(m<n\).

The questions are:
\begin{enumerate}
\item[(a)] What is \(f(3)\)?
\item[(b)] Is \(A_n\) bounded, equivalently does \(f(n)\) exist as a finite integer, for every integer \(n\geq 1\)? This question is attributed specifically to B. I. Plotkin.
\item[(c)] Is the growth of \(f(n)\) and \(g(n)\) polynomial, exponential, or intermediary? These growth categories and their comparison conventions are not defined in the supplied record; interpreting the question for \(f\) as a finite-valued function also depends on part (b).
\end{enumerate}
The supplied source reports, with attribution to N. Gupta and Yu. V. Kuz'min, that both \(f(n)-n\) and \(g(n)\) tend to infinity as \(n\to\infty\). This is retained as a supplied literature assertion, not as an independently verified result; its interpretation for \(f\) requires resolving the maximum convention.
\subsection{Short English statement}
For nilpotent groups of class exactly \(n\), determine how far nontrivial dimension subgroups can persist, compute \(f(3)\), and study the growth of this extremal index and the supplied upper-central-series quantity \(g(n)\). The coefficient ring, index conventions, existence of maxima, and meaning of the growth categories require review.
\subsection{Sources}
\begin{itemize}
\item Local UnsolvedMath registry, permanent ID \(701258\), exact source problem code \texttt{catalogue entry}, title \emph{Kourovka Notebook Problem 16.66}. The supplied registry labels the problem \emph{open}; that label is not independently verified here.
\item \emph{The Kourovka Notebook (Unsolved Problems in Group Theory)}, original problem 16.66, 16th issue (2006); supplied background also identifies the 21st-issue edition (2026), \texttt{arXiv:1401.0300v48}. The collection metadata names E. I. Khukhro and V. D. Mazurov as editors. Proposed by R. Mikhailov and I. B. S. Passi in 2006; part (b) is attributed to B. I. Plotkin.
\item N. Gupta and Yu. V. Kuz'min, \emph{J. Pure Appl. Algebra} \textbf{104} (1995), 191--197. This is the supplied reference for the assertion that \(f(n)-n\) and \(g(n)\) tend to infinity. No article title or source link is supplied.
\end{itemize}
\end{document}
